\section{Technical Overview}
\label{sec:overview}

This section is a self-contained tour: a reader who reads only it and the
introduction will understand how the argument works end to end, while the later
sections supply the definitions, lemmas, and proofs. The spine has four moves
and one unifying property. The arithmetization makes every constraint linear
except a handful of column products (\Cref{sec:overview-rep}); a binary-field
polynomial IOP discharges the linear constraints (\Cref{sec:overview-linear});
a single change of \emph{reading} lets the products ride that \emph{same} IOP
(\Cref{sec:overview-crux}); and a bit-sliced commitment settles everything with
one opening (\Cref{sec:overview-open}). What ties the four together is that
every map used to collapse a cell to the field is $\Ftwo$-linear
(\Cref{sec:overview-thread}).

\subsection{The representation: every primitive but $\AND$ is linear}
\label{sec:overview-rep}

A $\wbit$-bit word is the bit-polynomial $\bp{u} = \sum_b u_b X^b$ in
$\cellring = \Ftwo[X]/(X^{\wbit})$, the coefficient of $X^b$ being bit $b$.
$\XOR$ is then addition; a shift or rotation is multiplication by a monomial
$X^c$ in the shift quotient $\Fshift$ resp.\ rotation quotient $\Frot$, hence an
$\Ftwo$-linear map realised as a \emph{virtual} column reconstructed from its
source; and any $\Ftwo$-linear combination of these --- in particular the
$\Sigma/\sigma$ mixing functions, each an $\XOR$ of three rotations/shifts ---
is virtual too. None of these is committed or constrained.

Modular addition is the one arithmetic primitive. Taken in
$\Fshift = \Ftwo[X]/(X^{\wbit})$, where multiplication by $X$ is a left shift
and $X^{\wbit} = 0$ enacts the $\bmod\,2^{\wbit}$ wrap, a binary step
$\bp{s} \equiv \bp{a} + \bp{b} \pmod{2^{\wbit}}$ is pinned by the
\emph{Binius full-adder identity} (\Cref{lem:binius}): with the carry word
$\bp{c}$ defined by $\bp{c}[i] = (\bp{s}+\bp{a}+\bp{b})[i+1]$,
\begin{equation}
  \label{eq:ov-binius}
  \bp{s} = \bp{a} + \bp{b} + X\bp{c},
  \qquad
  (\bp{a} + X\bp{c}) \had (\bp{b} + X\bp{c}) = \bp{c} + X\bp{c}.
\end{equation}
The carry's lower $\wbit-1$ bits are the free virtual shift
$\SHR^{1}(\bp{s}+\bp{a}+\bp{b})$, so the only committed datum is the single
carry-out bit $\beta = \bp{c}[\wbit-1]$ that the shift cannot supply; the first
identity collapses to its bit-$0$ content $(\bp{s}+\bp{a}+\bp{b})[0]=0$, an
$\ideal{X}$-membership. The lone nonlinear primitive is bitwise $\AND$, the
coefficient-wise Hadamard product $\bp{u}\had\bp{v}$ (\Cref{lem:and}); $\Ch$ and
$\Maj$ each rewrite to a single such product. The arithmetization thus produces
an AIR in which \emph{every constraint is a linear ideal-membership relation
except a small fixed number of column Hadamard products} --- $K_H = 14$ of them
in the deployed layout ($2$ for the $\AND$s, $12$ for the addition carries). The
rest of the proof has exactly two jobs: the linear constraints, and the
Hadamard products.

\subsection{The linear part: ideal check, the reading $\psia$, sumcheck}
\label{sec:overview-linear}

The prover commits the primary witness columns natively in $\cellring$, with a
hash-based commitment whose $\Ftwo$-linear code \emph{does not widen} the cell
ring (\Cref{sec:commitment}). Because the trace already lives in
$\Ftwo[X]\subset\K[X]$, the prime-field projection that opens the general
Zinc\texttt{+} compiler is vacuous --- the first concrete saving of working in
characteristic two.

The linear constraints are ideal-membership assertions
$Q_\ell(T)\in\ideal{g_\ell}$ with generators $g_\ell\in\{0,\,X\}$. The
\emph{ideal check} (\Cref{sec:ideal-check}) batches all of them --- across
families and across all rows --- into combined polynomials $E_\ell\in\K[X]$
carrying the obligation $E_\ell\in\ideal{g_\ell}$. The \emph{evaluation reading}
$\psia$ then substitutes the cell variable $X=\alpha$ for a transcript-fresh
$\alpha\in\K$. Linear constraints survive this reading: $\psia$ is
$\Ftwo$-linear and evaluation is multiplicative, so
\[
  \psia\bigl(Q(T)\bigr) = Q\bigl(\psia(T)\bigr)
\]
exactly (\Cref{lem:linear-survives}). Hence a true linear identity among cells
becomes a true identity among the projected columns (completeness), while a
\emph{false} one survives only if a fixed nonzero degree-$<\wbit$ cell
polynomial vanishes at the random $\alpha$ --- probability $\le(\wbit-1)/|\K|$
(soundness, Schwartz--Zippel). A standard degree-$2$ sumcheck
(\Cref{sec:sumcheck}) then reduces the resulting zerocheck over the $2^{\nu}$
rows to per-column evaluation claims $\mle{\psia(w_g)}(r^*)$ at a single point
$r^*$.

\subsection{The crux: $\AND$ rides the same pipeline through a second reading $\psiz$}
\label{sec:overview-crux}

The Hadamard products cannot ride $\psia$, because evaluation does not respect a
coefficient-wise product:
\begin{equation}
  \label{eq:ov-scramble}
  \psia(\bp{u}\had\bp{v}) = \sum_b u_b v_b\,\alpha^{b}
  \;\neq\;
  \psia(\bp{u})\,\psia(\bp{v}) = \sum_{b,b'} u_b v_{b'}\,\alpha^{b+b'},
\end{equation}
the cross-terms $b\neq b'$ scrambling the bit-products
(\Cref{rem:and-not-survive}). The resolution is to read the \emph{same}
committed bits under a \emph{different} weighting. Identify the $\wbit$ bit
positions with an $\Ftwo$-affine subspace $H\subset\K$ ($\dim = \mu = \log_2\wbit$),
and read a cell's bits as the \emph{Lagrange} coefficients
$L(\bp{u}) = \sum_b u_b L_b$ of a univariate over $H$ rather than as monomial
coefficients. Then
\begin{equation}
  \label{eq:ov-had-ideal}
  \bp{u}\had\bp{v}=\bp{w}
  \quad\Longleftrightarrow\quad
  L(\bp{u})\,L(\bp{v}) - L(\bp{w}) \in \ideal{Z_H},
  \qquad Z_H = \textstyle\prod_{h\in H}(X-h)
\end{equation}
(\Cref{eq:had-ideal}): the product is correct on every bit exactly when this
polynomial vanishes on $H$. That is an \emph{ideal-membership claim} --- the
very shape the linear part already discharges, now with the extra generator
$Z_H$. So the nonlinearity rides the same machinery: the same ideal check (the
prover ships only the $\wbit-1$ off-subspace values of the combined product
polynomial, which vanish on $H$ for honest data), the same point $\alpha$ ---
now reading the cell in the \emph{Lagrange} basis,
\[
  \psiz(\bp{w}) = L(\bp{w})(\alpha) = \sum_b L_b(\alpha)\,w_b,
\]
rather than the monomial one --- and one extra degree-$3$ group in the
\emph{same} sumcheck, reduced to the \emph{same} point $r^*$
(\Cref{sec:hadamard}). The Lagrange reading keeps the bit-products $u_b v_b$
separated by position, so the verifier's closing check is a genuine
per-position product, not the scrambled cross-terms of \eqref{eq:ov-scramble}.
$\psia$ and $\psiz$ are not two projections: they are the one substitution
$X = \alpha$ with the bits read in the monomial resp.\ Lagrange basis (the
weights $\alpha^b = X^b|_\alpha$ and $L_b(\alpha) = L_b(X)|_\alpha$ are the same
point in two bases). The $\mu$ bit positions are folded once, by this reading,
and never become sumcheck variables.

\subsection{One opening answers both readings}
\label{sec:overview-open}

After the sumcheck and a multipoint-evaluation step (\Cref{sec:multipoint}),
both halves have reduced to multilinear-evaluation claims at one common point
$r_0$: the linear claims are $\psia$-images of the committed columns, the
discharge claims their $\psiz$-images. The commitment settles both from a single
object. Its opening binds, per column, the \emph{bit-slice fold}
\[
  a'_g = \sum_{b} a'_{g,b}\,X^b \in \K[X]_{<\wbit},
  \qquad
  a'_{g,b} = \sum_{x\in\{0,1\}^{\nu}} \eq_x(r_0)\,(w_g)_b(x),
\]
a degree-$<\wbit$ polynomial over $\K$ whose $b$-th coefficient is the $b$-th
bit slice of the column, folded along $\eq(\cdot;r_0)$. Because every reading is
the inner product of the coefficient vector with its weights, \emph{any}
projected evaluation at $r_0$ is a single length-$\wbit$ read-off of $a'_g$:
\[
  \psia(a'_g) = \mle{\psia(w_g)}(r_0)\ \ (\text{linear claim}),
  \qquad
  \psiz(a'_g) = \mle{\psiz(w_g)}(r_0)\ \ (\text{discharge claim}).
\]
One bound object answers both readings, so the nonlinear constraints cost
\emph{no second opening and no second point} (\Cref{sec:commitment}). The
opening itself is a flat Brakedown/Ligero-style proximity test
\cite{brakedown} over the $\Ftwo$-RAA code, with knowledge-soundness error
$\epsilon_{\mathrm{prox}}(t,\delta) + O(1/|\K|)$ dominated by its $t$ column
queries.

\subsection{The thread, the guarantee, and the scope}
\label{sec:overview-thread}

One property runs through all four moves: every reading is an $\Ftwo$-linear map
$\cellring\to\K$, $\sum_b w_b X^b \mapsto \sum_b w_b\,\xi_b$, and the
commitment's code is $\Ftwo$-linear and does not widen $\cellring$. Hence
readings commute through encoding, $\XOR$ of columns commutes with the reading,
shift and rotation virtuals commute with both, the random reading may be applied
\emph{after} committing, and one committed object answers many functionals
(\Cref{rem:linearity-thread}). That single fact is what makes $\XOR$ free,
shifts virtual, and the $\psia/\psiz$ pair share an opening.

Composed under Fiat--Shamir, the argument is a knowledge-sound SNARK for the
SHA-256 chained-compression relation, with error
\[
  \epsilon \;\le\; \epsilon_{\mathrm{prox}}(t,\delta)
  + O\!\Bigl(\tfrac{N_{\mathrm{con}}\,\wbit + K_H + \nu}{|\K|}\Bigr)
\]
(\Cref{thm:e2e}), the proximity term dominating. Two honest caveats remain
(\Cref{sec:caveats}): the construction is not zero-knowledge, and the
modular-addition carries are presently bound by a trusted reading rather than to
the commitment --- the $\AND$ relations and the entire linear part are
unconditional, and closing the adder gap is a localized, recorded follow-up.
